% Resumo da prévia do TCC2; atualizar após incorporar os resultados experimentais.
\begin{resumoutfpr}
Este trabalho apresenta o desenvolvimento de um jogo tridimensional multijogador competitivo de tiro em primeira pessoa, com combate em arena e elementos de plataforma. O objetivo foi integrar movimentação, combate, recursos coletáveis, modos de jogo e comunicação em rede em um produto desenvolvido com a Godot Engine 4.5 e GDScript. A implementação foi organizada em dois projetos, Client e Server, e em um diretório de definições compartilhadas. O Client concentra a entrada e a apresentação, enquanto um servidor dedicado mantém as regras de dano, recursos, pontuação e encerramento das partidas, utilizando posições informadas pelos jogadores. A versão descrita reúne salto duplo, plataformas de impulso, três armas, granadas, pickups e efeitos temporários, com partidas em equipes nos modos Team Deathmatch e King of the Hill no mapa Highrise. A sincronização utiliza chamadas remotas de procedimento, snapshots, interpolação de estados remotos e consulta ao histórico para validar disparos. Também foram implementados coleta de métricas, registro em CSV e uma dashboard para acompanhamento da execução. O relatório apresenta a organização por responsabilidades, as decisões de sincronização e o funcionamento dos sistemas desenvolvidos, acompanhados de suas limitações. Os registros dos testes integrados e os resultados quantitativos permanecem como evidências a incorporar à versão final.
\end{resumoutfpr}
